%% main.tex —— 东北财经大学研究生学位论文 LaTeX 模板（示例文档）
%% 编译：XeLaTeX 连续编译 3 次，或直接运行 ./compile.sh
%%
%% 说明：本文件是一份「示例论文」，题目、作者、导师、数据等全部为虚构，
%%       仅用于演示模板排版。请将各处替换为你自己的真实内容。

%% ============================================================
%% 学位类型：只改这一处即可在 学硕 / 专硕 / 博士 之间切换
%% ============================================================
\documentclass[master]{dufethesis}
% master     —— 学术硕士（默认）：封面中间为「一级学科名称 / 二级学科名称」
% profmaster —— 专业硕士：封面中间为「学科专业名称 / 所在学院名称」
% doctor     —— 博士：封面顶部显示「博士学位论文」

%% ============================================================
%% 排版基础设施（一般无需改动）
%% ============================================================
%% 自绘 algorithm / breakablealgorithm 浮动环境（兼容缺失 algorithm.sty 的精简发行版）
\newcounter{algctr}[chapter]
\renewcommand{\thealgctr}{\thechapter.\arabic{algctr}}
\newenvironment{algorithm}[1][H]{%
  \begin{figure}[#1]\refstepcounter{algctr}\centering\small%
  \begin{minipage}{0.94\textwidth}\hrule height 0.6pt\vspace{0.4em}%
}{%
  \vspace{0.4em}\hrule height 0.4pt\end{minipage}\end{figure}%
}
\newenvironment{breakablealgorithm}{%
  \dufealgneedspace
  \refstepcounter{algctr}%
  \enlargethispage{1.2\baselineskip}%
  \par\addvspace{0.6em}\begingroup\small\setstretch{1.3}\setlength{\parskip}{0pt}%
  \noindent\rule{\textwidth}{0.6pt}\par%
}{%
  \par\vspace{0.1em}\noindent\rule{\textwidth}{0.4pt}%
  \endgroup\addvspace{0.6em}%
}
\newcommand{\dufealgneedspace}{%
  \par\begingroup
  \dimen0=\pagegoal \advance\dimen0 by -\pagetotal
  \ifdim\dimen0<10\baselineskip\pagebreak\fi
  \endgroup
}
\newcommand{\algcaption}[1]{%
  \noindent\hbox to\textwidth{\hfil\textbf{算法~\thealgctr~~#1}\hfil}\par\nopagebreak[4]%
  \noindent\rule{\textwidth}{0.3pt}\par\nopagebreak[4]\vspace{0.45em}\nopagebreak[4]%
}

%% 图 / 表「资料来源」标注（小五宋体，置于图表左下方）
\newcommand{\dufesource}[1]{\par\vspace{3pt}\noindent{\xiaowu\songti\parbox{\textwidth}{#1}}}

%% 正文作者-年份引用：用 \citep{键} 标注；参考文献列表仍按学校要求用 [N] 编号。
%% 维护引用库：\addcite{键}{作者, 年份}，需与下方 references 列表一一对应。
\makeatletter
\newcommand{\addcite}[2]{\@namedef{cite@#1}{#2}}
\newcommand{\thesiscite}[1]{\@ifundefined{cite@#1}{\textbf{[?#1]}}{\@nameuse{cite@#1}}}
\newcommand{\citep}[1]{（\dufe@citep#1,\@nil）}
\def\dufe@citep#1,#2\@nil{%
  \thesiscite{#1}%
  \if\relax\detokenize{#2}\relax\else；\dufe@citep#2\@nil\fi}
\makeatother

%% 允许 TeX 额外拉伸段落间距，缓解 CJK+英文混排造成的少量 overfull
\emergencystretch=2em
\tolerance=2000

%% ============================================================
%% 引用库（示例）—— 每条对应下方 references 列表中的一项
%% ============================================================
\addcite{xugc2023digital}{许宪春和张美慧, 2023}
\addcite{liwei2022region}{李伟, 2022}
\addcite{smith2020measure}{Smith and Jones, 2020}
\addcite{chen2019cluster}{Chen et al., 2019}

%% ============================================================
%% 论文信息（示例：请全部替换为你自己的信息）
%% ============================================================
\schoolcode{10173}                 % 学校编码（固定）
\classno{}                         % 分类号（留空，提交时填写）
\secretlevel{}                     % 密级（公开论文留空）
\studentid{2024110001}             % 学号
\udc{}                             % UDC（留空，提交时填写）

\title{数字经济发展水平测度与区域差异研究}
\entitle{Measurement and Regional Disparity of Digital Economy Development}
\author{张三}
\supervisor{李四}
\supervisortitle{教授}

%% 学硕 / 博士 显示以下两项：
\major{应用经济学}                  % 一级学科
\submajor{统计学}                   % 二级学科

%% 专硕 显示以下三项（选择 profmaster 时生效）：
\proftype{应用统计}                 % 学科专业名称
\college{统计与数学学院}             % 所在学院名称
\submitdate{2025年5月}              % 论文提交日期

\defenddate{2025年5月}             % 论文答辩时间

\begin{document}

%% ===== 封面 =====
\makecover

%% ===== 前置部分（罗马数字页码）=====
\frontmatter

%% ----- 中文摘要 -----
\begin{cabstract}
数字经济已成为推动高质量发展的重要引擎，科学测度其发展水平并刻画区域差异，对优化资源配置具有重要意义。本文以虚构的省际面板数据为例，从数字基础设施、数字产业化、产业数字化与数字治理四个维度构建数字经济发展水平综合指标，采用熵值法确定权重，并结合泰尔指数对区域差异进行分解。

研究发现，样本期内数字经济发展水平总体呈上升趋势，但区域间不平衡特征明显，差异主要来源于组间差距。基于上述结果，本文从基础设施补短板、产业协同与制度供给等方面提出对策建议。

需要特别说明的是：本摘要与全文内容均为模板演示用的虚构示例，不代表任何真实研究结论；请在使用本模板时替换为实际内容。摘要应具有独立性和自含性，硕士论文一般 1500 字左右，博士论文 3000 字左右，且不出现图、表与数学公式。
\end{cabstract}
\ckeywords{数字经济，综合测度，区域差异，熵值法，泰尔指数}

%% ----- 英文摘要 -----
\begin{eabstract}
The digital economy has become a key engine of high-quality development, and measuring its level while characterizing regional disparities is of great significance for optimizing resource allocation. Using fictitious provincial panel data as an example, this paper constructs a composite index of digital-economy development from four dimensions---digital infrastructure, digital industrialization, industrial digitalization, and digital governance---determines the weights via the entropy method, and decomposes regional disparities using the Theil index.

The results show an overall upward trend over the sample period, alongside pronounced regional imbalance driven mainly by between-group gaps. Corresponding policy implications are then discussed.

Please note that the abstract and the entire document are fictitious examples for template demonstration only and do not represent any real research findings. Replace them with your own content when using this template.
\end{eabstract}
\ekeywords{Digital Economy, Composite Measurement, Regional Disparity, Entropy Method, Theil Index}

%% ----- 目录 -----
\tableofcontents

%% 图表目录（博士论文一般需要，按需取消注释）
% \listoffigures
% \listoftables

%% ===== 正文部分（阿拉伯数字页码）=====
\mainmatter

\chapter{绪论}
\label{chap:intro}

\section{研究背景与意义}

近年来，数字经济在国民经济中的占比持续提升，正深刻改变要素配置方式与产业组织形态。然而，由于资源禀赋与发展阶段不同，各地区数字经济发展呈现明显的不平衡特征\citep{xugc2023digital}。如何科学测度数字经济发展水平、客观刻画其区域差异，已成为经济统计领域的重要议题\citep{liwei2022region}。

本文的研究意义主要体现在两个方面。其一，理论上，构建一套多维度、可比较的数字经济发展水平测度框架，有助于丰富数字经济量化研究\footnote{本文为模板演示文档，全部数据与结论均为虚构，仅用于展示排版效果。}。其二，实践上，区域差异的来源分解可为差异化政策制定提供参考。

\section{研究内容与研究方法}

围绕“测度—比较—解释”这一主线，本文拟回答以下问题：

\begin{enumerate}[label=（\arabic*）]
    \item 如何从多个维度构建数字经济发展水平的综合指标？
    \item 样本期内数字经济发展水平的总体趋势如何？
    \item 区域差异主要来源于组内差距还是组间差距？
\end{enumerate}

为此，本文综合运用以下方法：采用熵值法客观确定指标权重，避免主观赋权的随意性；采用泰尔指数对区域差异进行可加分解；并辅以聚类分析对地区进行分组比较\citep{chen2019cluster}。

\section{研究框架}

本文共分五章。第~\ref{chap:method}~章给出指标体系与测度方法，总体技术路线见图~\ref{fig:roadmap}；第~\ref{chap:empirical}~章在虚构数据上展开实证分析；末章总结结论并讨论展望。
     % 第1章 绪论
\chapter{文献综述与理论基础}
\label{chap:litreview}

\section{数字经济测度研究综述}

数字经济测度大致沿两条路径展开：一是基于增加值核算的规模测算，二是基于综合指标体系的水平评价\citep{xugc2023digital}。前者口径清晰但对统计数据要求较高；后者灵活性强、便于跨地区横向比较，但权重确定方法对结果影响较大\citep{liwei2022region}。本文采用综合指标体系路径，并以客观赋权缓解主观随意性。

\section{区域差异的测度与分解}

刻画区域差异的常用工具包括基尼系数、变异系数与泰尔指数。其中泰尔指数因满足可加分解性而被广泛采用：总差异可分解为组内差异与组间差异之和。设地区 $i$ 的数字经济发展水平占总体的比重为 $s_i$，地区个数为 $n$，则泰尔指数定义为

\begin{equation}
    T = \sum_{i=1}^{n} s_i \ln\!\frac{s_i}{1/n}.
    \label{eq:theil}
\end{equation}

式~\eqref{eq:theil} 取值越大，表示区域间发展越不均衡。当样本被划分为若干组时，$T$ 可进一步分解为组间项与组内项之和，从而识别差异的主要来源。

\section{理论基础}

数字经济的区域非均衡可由要素集聚与空间溢出理论加以解释：数字基础设施与人才在中心地区集聚，形成规模报酬递增，进而拉大区域差距\citep{chen2019cluster}。这一理论视角为后文的差异分解提供了解释框架。
     % 第2章 文献综述与理论基础
\chapter{研究设计与测度方法}
\label{chap:method}

\section{指标体系构建}

本文从多个维度构建数字经济发展水平综合评价体系，下文先说明维度划分，再说明指标选取与标准化处理。

\subsection{维度划分}

综合评价体系包含数字基础设施、数字产业化、产业数字化与数字治理四个一级维度，分别刻画数字经济发展的支撑条件、产业形态与治理环境。

\subsection{指标选取与标准化}

各一级维度下进一步选取若干可量化的基础指标，并经极差法标准化后进入综合测度，标准化方法见附录~\ref{chap:appendix}。

\section{基于熵值法的综合测度}

为客观确定权重，本文采用熵值法：信息熵越小的指标，离散程度越大、所提供的信息越多，应赋予越大的权重。综合测度的主要步骤如算法~\ref{alg:entropy} 所示。

\begin{breakablealgorithm}
\algcaption{基于熵值法的综合测度}
\label{alg:entropy}
\noindent\textbf{输入}：标准化指标矩阵 $X=(x_{ij})_{m\times n}$（$m$ 个地区、$n$ 项指标）\\
\textbf{输出}：各地区综合得分 $\{S_i\}$
\begin{enumerate}[label=\arabic*., itemsep=1.6pt, topsep=2pt, parsep=0pt, leftmargin=2.0em]
    \item 计算指标 $j$ 下地区 $i$ 的比重 $p_{ij}=x_{ij}/\sum_{i} x_{ij}$。
    \item 计算指标 $j$ 的信息熵 $e_j=-k\sum_{i} p_{ij}\ln p_{ij}$，其中 $k=1/\ln m$。
    \item 计算差异系数 $g_j=1-e_j$，归一化得到权重 $w_j=g_j/\sum_j g_j$。
    \item 加权求和得到综合得分 $S_i=\sum_j w_j x_{ij}$。
    \item \textbf{return} $\{S_i\}$。
\end{enumerate}
\end{breakablealgorithm}

\section{技术路线}

本文整体研究流程如图~\ref{fig:roadmap} 所示：从指标选取与数据标准化出发，经熵值法测度得到综合得分，再进行区域差异分解与聚类比较，最终形成结论与建议。

\begin{figure}[!htbp]
    \centering
    \includegraphics[width=0.8\textwidth]{figures/demo.png}
    \caption{本文研究技术路线图}
    \label{fig:roadmap}
\end{figure}
     % 第3章 研究设计与测度方法
\chapter{实证分析}
\label{chap:empirical}

\section{数据来源与说明}

本文使用虚构的 2018—2022 年省际面板数据进行演示，所有数值均为模拟生成，不代表真实情况。为便于比较，将样本地区划分为东部、中部与西部三组。

\section{综合测度结果}

按第~\ref{chap:method}~章方法测算，部分地区数字经济发展水平综合得分如表~\ref{tab:score} 所示。可以看到，东部地区得分整体高于中西部，且呈逐年上升趋势。

\begin{table}[!htbp]
    \centering
    \caption{部分地区数字经济发展水平综合得分（示例）}
    \label{tab:score}
    \wuhao\songti
    \begin{tabular}{lccc}
        \toprule
        地区 & 2018 年 & 2020 年 & 2022 年 \\
        \midrule
        A 省 & 0.412 & 0.498 & 0.576 \\
        B 省 & 0.305 & 0.361 & 0.423 \\
        C 省 & 0.241 & 0.287 & 0.339 \\
        \bottomrule
    \end{tabular}

    \dufesource{资料来源：本文基于虚构数据测算整理。}
\end{table}

\section{区域差异分解}

为识别区域差异的来源，本文先测算总体泰尔指数，再将其分解为组间与组内两部分。

\subsection{泰尔指数测算}

按式~\eqref{eq:theil} 计算总体泰尔指数，结果见表~\ref{tab:theil}。

\begin{table}[!htbp]
    \centering
    \caption{泰尔指数分解结果（示例）}
    \label{tab:theil}
    \wuhao\songti
    \begin{tabular}{clcc}
        \toprule
        年份 & 差异来源 & 贡献值 & 贡献率 \\
        \midrule
        \multirow{2}{*}{2018} & 组间差异 & 0.061 & 64\% \\
                              & 组内差异 & 0.034 & 36\% \\
        \midrule
        \multirow{2}{*}{2022} & 组间差异 & 0.078 & 69\% \\
                              & 组内差异 & 0.035 & 31\% \\
        \bottomrule
    \end{tabular}

    \dufesource{资料来源：本文基于虚构数据测算整理。}
\end{table}

\subsection{组间与组内分解}

将总差异分解后可见，区域差异主要来源于组间差距。综合得分的区域分布如图~\ref{fig:dist} 所示，直观印证了上述格局。

\begin{figure}[!htbp]
    \centering
    \includegraphics[width=0.8\textwidth]{figures/demo.png}
    \caption{综合得分的区域分布（示例）}
    \label{fig:dist}
\end{figure}
     % 第4章 实证分析
\chapter{研究结论与展望}
\label{chap:conclusion}

\section{主要结论}

基于前文的测度与分解，本文得到以下主要结论：

\begin{enumerate}[label=（\arabic*）]
    \item 样本期内数字经济发展水平总体呈上升趋势，但地区间绝对差距仍在扩大；
    \item 区域差异主要来源于组间差距，东部与中西部之间的差距是总差异的主要构成；
    \item 数字基础设施与数字产业化维度对综合得分的贡献相对更大。
\end{enumerate}

\section{政策建议}

针对上述结论，本文提出三点建议：加大中西部数字基础设施投入以补齐短板；推动区域间产业协同与数据要素流动；完善数字治理的制度供给，营造良好的数字经济发展环境。

\section{不足与展望}

受演示数据所限，本文结论仅用于说明方法流程。未来可在真实数据基础上，引入空间计量模型刻画区域间溢出效应，并延长样本期以考察数字经济发展的动态演变趋势。
     % 第5章 研究结论与展望

%% ===== 附录（属于正文编号体系：附录 A、附录 B……，故置于 \backmatter 之前）=====
\appendix
\clearpage
\refstepcounter{chapter}
\phantomsection
\addcontentsline{toc}{chapter}{附录 \thechapter\hspace{0.5em}补充结果}
\chapter*{附录 \thechapter\hspace{1em}补充结果}
\markboth{附录 \thechapter\hspace{1em}补充结果}{}
\label{chap:appendix}

本附录给出补充的指标说明与标准化公式，用于演示附录中表、公式的编号方式（表~\ref{tab:appendix_indicators}、式~\eqref{eq:appendix_norm}）。

\section{指标说明}

表~\ref{tab:appendix_indicators} 列出数字经济发展水平评价的部分基础指标及其单位。

\begin{table}[!htbp]
    \centering
    \caption{数字经济发展水平评价部分基础指标（示例）}
    \label{tab:appendix_indicators}
    \wuhao\songti
    \begin{tabular}{llc}
        \toprule
        维度 & 指标 & 单位 \\
        \midrule
        数字基础设施 & 每百人互联网宽带接入端口数 & 个 \\
        数字产业化   & 信息传输、软件业增加值占比 & \% \\
        产业数字化   & 电子商务销售额 & 亿元 \\
        \bottomrule
    \end{tabular}

    \dufesource{资料来源：本文整理（示例）。}
\end{table}

\section{标准化方法}

各基础指标采用极差法进行标准化，正向指标的标准化公式为

\begin{equation}
    x'_{ij} = \frac{x_{ij}-\min_i x_{ij}}{\max_i x_{ij}-\min_i x_{ij}}.
    \label{eq:appendix_norm}
\end{equation}


%% ===== 后置部分（参考文献、后记等，章节不再编号）=====
\backmatter

%% ----- 在学期间发表的科研成果（仅博士论文需要，按需取消注释）-----
% \begin{publications}
% \textbf{一、发表的学术论文}
%
% 1. 作者姓名. 论文标题[J]. 期刊名称, 2024, (1): 1-10.
% \end{publications}

%% ----- 参考文献（与上方 \addcite 一一对应）-----
\begin{references}
    \item 许宪春，张美慧. 中国数字经济规模测算研究与展望[M]. 北京：中国统计出版社，2023. 35-58.
    \item 李伟. 区域数字经济发展水平测度与差异分解[J]. 统计研究，2022，39（6）：45-58.
    \item 王明亮. 关于中国学术期刊标准化数据库系统工程的进展[EB/OL]. http://www.example.org/digital.html，2023-06-01/2023-09-10.
    \item 张华. 数字经济与区域协调发展[N]. 经济日报，2023-03-12（10）.
    \item {\enref Smith, J., and Jones, K., 2020, "Measuring the Digital Economy", Journal of Economic Statistics, Vol.38, No.2, PP120-145.}
    \item {\enref Chen, L., Wang, M., and Liu, P., 2019, Cluster Analysis for Regional Development, Springer Press.}
    % 添加更多参考文献……硕士论文不少于 40 项，博士论文不少于 100 项
\end{references}

%% ----- 后记 -----
\begin{acknowledgements}
本论文是在导师李四教授的悉心指导下完成的。从选题、框架到反复修改，导师严谨的治学态度和诲人不倦的精神令我受益匪浅，在此谨致以最诚挚的谢意。

感谢学院各位老师在课程学习与开题过程中给予的宝贵建议，感谢同窗好友的讨论与帮助，也感谢家人一路以来的理解与支持。

（本后记为模板示例文字；后记语言宜简明诚恳，字数 400 字左右。）
\end{acknowledgements}
\ackauthor{张三}{2025年5月}

%% ===== 原创性声明与授权书 =====
\makedeclaration

\end{document}
